\documentclass[10pt,a4paper]{amsart}
\usepackage{amsmath,amssymb,amsthm,graphicx}
\usepackage{verbatim}
\newenvironment{reference}{\comment}{\endcomment}
%\newenvironment{reference}{}{}
\newenvironment{slogan}{\comment}{\endcomment}
\newenvironment{history}{\comment}{\endcomment}

% For commutative diagrams we use Xy-pic
\usepackage[all]{xy}

% We use 2cell for 2-commutative diagrams.
\xyoption{2cell}
% Theorem environments.
%
\theoremstyle{plain}
\newtheorem{theorem}[subsection]{Theorem}
\newtheorem{proposition}[subsection]{Proposition}
\newtheorem{lemma}[subsection]{Lemma}

\theoremstyle{definition}
\newtheorem{definition}[subsection]{Definition}
\newtheorem{example}[subsection]{Example}
\newtheorem{exercise}[subsection]{Exercise}
\newtheorem{situation}[subsection]{Situation}

\theoremstyle{remark}
\newtheorem{remark}[subsection]{Remark}
\newtheorem{remarks}[subsection]{Remarks}

\numberwithin{equation}{subsection}

% Macros
%
\def\lim{\mathop{\mathrm{lim}}\nolimits}
\def\colim{\mathop{\mathrm{colim}}\nolimits}
\def\Spec{\mathop{\mathrm{Spec}}}
\def\Hom{\mathop{\mathrm{Hom}}\nolimits}
\def\Ext{\mathop{\mathrm{Ext}}\nolimits}
\def\SheafHom{\mathop{\mathcal{H}\!\mathit{om}}\nolimits}
\def\SheafExt{\mathop{\mathcal{E}\!\mathit{xt}}\nolimits}
\def\Sch{\mathit{Sch}}
\def\Mor{\mathop{\mathrm{Mor}}\nolimits}
\def\Ob{\mathop{\mathrm{Ob}}\nolimits}
\def\Sh{\mathop{\mathit{Sh}}\nolimits}
\def\NL{\mathop{N\!L}\nolimits}
\def\CH{\mathop{\mathrm{CH}}\nolimits}
\def\proetale{{pro\text{-}\acute{e}tale}}
\def\etale{{\acute{e}tale}}
\def\QCoh{\mathit{QCoh}}
\def\Ker{\mathop{\mathrm{Ker}}}
\def\Im{\mathop{\mathrm{Im}}}
\def\Coker{\mathop{\mathrm{Coker}}}
\def\Coim{\mathop{\mathrm{Coim}}}

% Boxtimes
%
\DeclareMathSymbol{\boxtimes}{\mathbin}{AMSa}{"02}

%
% Macros for moduli stacks/spaces
%
\def\QCohstack{\mathcal{QC}\!\mathit{oh}}
\def\Cohstack{\mathcal{C}\!\mathit{oh}}
\def\Spacesstack{\mathcal{S}\!\mathit{paces}}
\def\Quotfunctor{\mathrm{Quot}}
\def\Hilbfunctor{\mathrm{Hilb}}
\def\Curvesstack{\mathcal{C}\!\mathit{urves}}
\def\Polarizedstack{\mathcal{P}\!\mathit{olarized}}
\def\Complexesstack{\mathcal{C}\!\mathit{omplexes}}
% \Pic is the operator that assigns to X its picard group, usage \Pic(X)
% \Picardstack_{X/B} denotes the Picard stack of X over B
% \Picardfunctor_{X/B} denotes the Picard functor of X over B
\def\Pic{\mathop{\mathrm{Pic}}\nolimits}
\def\Picardstack{\mathcal{P}\!\mathit{ic}}
\def\Picardfunctor{\mathrm{Pic}}
\def\Deformationcategory{\mathcal{D}\!\mathit{ef}}
\usepackage{cleveref}
\crefname{originalequation}{original equation}{original equations}
\Crefname{originalequation}{Original equation}{Original equations}
\makeatletter
\newlabel{equation-bijective-orbits}{{06T6}{source}{}{sourcepointer-1}{}}
\newlabel{equation-bijective-orbits@cref}{{[originalequation][0][]06T6}{[1][0][]source}{}{}{}}
\makeatother
\crefname{originalreference}{original reference}{original references}
\Crefname{originalreference}{Original reference}{Original references}
\makeatletter
\newlabel{example-prorepresentable-deformation-functor}{{06JB}{source}{}{sourcepointer-2}{}}
\newlabel{example-prorepresentable-deformation-functor@cref}{{[originalreference][0][]06JB}{[1][0][]source}{}{}{}}
\makeatother
\crefname{originalreference}{original reference}{original references}
\Crefname{originalreference}{Original reference}{Original references}
\makeatletter
\newlabel{example-tangent-space-map-prorepresentable-functor}{{06ST}{source}{}{sourcepointer-3}{}}
\newlabel{example-tangent-space-map-prorepresentable-functor@cref}{{[originalreference][0][]06ST}{[1][0][]source}{}{}{}}
\makeatother
\crefname{originalreference}{original reference}{original references}
\Crefname{originalreference}{Original reference}{Original references}
\makeatletter
\newlabel{example-tangent-space-prorepresentable-functor}{{06IC}{source}{}{sourcepointer-4}{}}
\newlabel{example-tangent-space-prorepresentable-functor@cref}{{[originalreference][0][]06IC}{[1][0][]source}{}{}{}}
\makeatother
\crefname{originalreference}{original reference}{original references}
\Crefname{originalreference}{Original reference}{Original references}
\makeatletter
\newlabel{lemma-RS-implies-S1-S2}{{06J7}{source}{}{sourcepointer-5}{}}
\newlabel{lemma-RS-implies-S1-S2@cref}{{[originalreference][0][]06J7}{[1][0][]source}{}{}{}}
\makeatother
\crefname{originalreference}{original reference}{original references}
\Crefname{originalreference}{Original reference}{Original references}
\makeatletter
\newlabel{lemma-action-linear}{{06SU}{source}{}{sourcepointer-6}{}}
\newlabel{lemma-action-linear@cref}{{[originalreference][0][]06SU}{[1][0][]source}{}{}{}}
\makeatother
\crefname{originalreference}{original reference}{original references}
\Crefname{originalreference}{Original reference}{Original references}
\makeatletter
\newlabel{lemma-minimal-smooth-morphism-functors}{{06JL}{source}{}{sourcepointer-7}{}}
\newlabel{lemma-minimal-smooth-morphism-functors@cref}{{[originalreference][0][]06JL}{[1][0][]source}{}{}{}}
\makeatother
\crefname{originalreference}{original reference}{original references}
\Crefname{originalreference}{Original reference}{Original references}
\makeatletter
\newlabel{theorem-miniversal-object-existence}{{06IX}{source}{}{sourcepointer-8}{}}
\newlabel{theorem-miniversal-object-existence@cref}{{[originalreference][0][]06IX}{[1][0][]source}{}{}{}}
\makeatother
\title{Criterion for prorepresentability of a functor on Artinian algebras}
\author{The Stacks Project Authors}
\date{stacks-project:formal-defos.tex:local-snapshot-20261012; GNU FDL 1.2 or later}
\begin{document}
\maketitle
\noindent\textit{Editorial note.} All delivered mathematical text is the author's own, copied from the pinned Stacks Project chapter file formal-defos.tex; no statement or proof step was rewritten, reordered or removed.\par\medskip
\noindent\textit{Editorial note.} The counted proof is the author's complete argument: the three verifications for a prorepresentable functor, the passage to a minimal versal formal object via the deformation functor axioms, and the closing orbit argument upgrading the transformation to an isomorphism.\par\medskip
\noindent\textit{Editorial note.} The label theorem-Schlessinger-prorepresentability names the counted statement, and the proof environment immediately following it in the source is delivered verbatim as its proof.\par\medskip
\noindent\textit{Editorial note.} The counted statement is carried by the author's own theorem environment in the chapter source.\par\medskip
\noindent\textit{Editorial note.} Pre-open gate: python operations/agent-authoring/tools/precard\_gate.py --source formal-defos.tex --label theorem-Schlessinger-prorepresentability --start 5738 --end 5775 --local-id 28027 returned RESULT=OK (L28027 free, occupancy FREE, public interval clear). Candidate taken from sources/stacks-formal-defos/candidates-halfA.json.\par\medskip
\noindent\textit{Editorial note.} The statement and proof are copied from the pinned chapter file; the cited inputs are kept as external pointers with their canonical Stacks tags.\par\medskip
\noindent\textit{Allowed standard prerequisites (editorial).}\begin{itemize}
\item Prorepresentable functors are deformation functors with finite dimensional tangent space. Same chapter, the examples establishing that a prorepresentable functor is a deformation functor, that its tangent space is finite dimensional and that the comparison with derivations is as displayed. Supplies the three verifications in the first direction of the theorem.
\item (RS) implies (S1) and (S2) for deformation functors. Same chapter, the lemma deriving (S1) and (S2) from the deformation functor axioms. Supplies (S1) and (S2) for the functor under the hypotheses of the converse.
\item Existence of a minimal versal formal object satisfying the orbit condition. Same chapter, the theorem producing a minimal versal formal object whose differential has bijective orbits. Supplies the formal object xi lying over R used in the converse.
\item Minimality criterion and the action by derivations. Same chapter, the characterization of minimality for smooth morphisms of functors and the lemma describing the action of the derivations of the residue field by addition. Supply the conclusion that the map on tangent spaces is bijective and hence that the transformation is an isomorphism.
\end{itemize}\medskip
\ifdefined\hypertarget\hypertarget{sourcepointer-1}{}\fi
\noindent\textit{Source reference.} Original equation 06T6: sources/stacks-formal-defos/pinned-original-source.tex, source lines 4623--4623. Prerequisite pointer invoked at proof line(s) 5764; it names a standard fact of the same chapter that the counted argument uses, and every step of the counted proof built on it is retained verbatim in the delivered body.\par\medskip
\ifdefined\hypertarget\hypertarget{sourcepointer-2}{}\fi
\noindent\textit{Source reference.} Original reference 06JB: sources/stacks-formal-defos/pinned-original-source.tex, source lines 5328--5328. Prerequisite pointer invoked at proof line(s) 5750; it names a standard fact of the same chapter that the counted argument uses, and every step of the counted proof built on it is retained verbatim in the delivered body.\par\medskip
\ifdefined\hypertarget\hypertarget{sourcepointer-3}{}\fi
\noindent\textit{Source reference.} Original reference 06ST: sources/stacks-formal-defos/pinned-original-source.tex, source lines 3606--3606. Prerequisite pointer invoked at proof line(s) 5755; it names a standard fact of the same chapter that the counted argument uses, and every step of the counted proof built on it is retained verbatim in the delivered body.\par\medskip
\ifdefined\hypertarget\hypertarget{sourcepointer-4}{}\fi
\noindent\textit{Source reference.} Original reference 06IC: sources/stacks-formal-defos/pinned-original-source.tex, source lines 3545--3545. Prerequisite pointer invoked at proof line(s) 5752; it names a standard fact of the same chapter that the counted argument uses, and every step of the counted proof built on it is retained verbatim in the delivered body.\par\medskip
\ifdefined\hypertarget\hypertarget{sourcepointer-5}{}\fi
\noindent\textit{Source reference.} Original reference 06J7: sources/stacks-formal-defos/pinned-original-source.tex, source lines 5164--5164. Prerequisite pointer invoked at proof line(s) 5760; it names a standard fact of the same chapter that the counted argument uses, and every step of the counted proof built on it is retained verbatim in the delivered body.\par\medskip
\ifdefined\hypertarget\hypertarget{sourcepointer-6}{}\fi
\noindent\textit{Source reference.} Original reference 06SU: sources/stacks-formal-defos/pinned-original-source.tex, source lines 3811--3811. Prerequisite pointer invoked at proof line(s) 5771; it names a standard fact of the same chapter that the counted argument uses, and every step of the counted proof built on it is retained verbatim in the delivered body.\par\medskip
\ifdefined\hypertarget\hypertarget{sourcepointer-7}{}\fi
\noindent\textit{Source reference.} Original reference 06JL: sources/stacks-formal-defos/pinned-original-source.tex, source lines 5695--5695. Prerequisite pointer invoked at proof line(s) 5774; it names a standard fact of the same chapter that the counted argument uses, and every step of the counted proof built on it is retained verbatim in the delivered body.\par\medskip
\ifdefined\hypertarget\hypertarget{sourcepointer-8}{}\fi
\noindent\textit{Source reference.} Original reference 06IX: sources/stacks-formal-defos/pinned-original-source.tex, source lines 4837--4837. Prerequisite pointer invoked at proof line(s) 5762; it names a standard fact of the same chapter that the counted argument uses, and every step of the counted proof built on it is retained verbatim in the delivered body.\par\medskip
\begin{theorem}
\label{theorem-Schlessinger-prorepresentability}
Let $F: \mathcal{C}_\Lambda \to \textit{Sets}$ be a functor.
Then $F$ is prorepresentable if and only if
(a) $F$ is a deformation functor,
(b) $\dim_k TF$ is finite, and (c) $\gamma : \text{Der}_\Lambda(k, k) \to TF$
is injective.
\end{theorem}

\begin{proof}
Assume $F$ is prorepresentable by $R \in \widehat{\mathcal{C}}_\Lambda$.
We see $F$ is a deformation functor by
Example \ref{example-prorepresentable-deformation-functor}.
We see $\dim_k TF$ is finite by
Example \ref{example-tangent-space-prorepresentable-functor}.
Finally, $\text{Der}_\Lambda(k, k) \to TF$ is identified with
$\text{Der}_\Lambda(k, k) \to \text{Der}_\Lambda(R, k)$ by
Example \ref{example-tangent-space-map-prorepresentable-functor}
which is injective because $R \to k$ is surjective.

\medskip\noindent
Conversely, assume (a), (b), and (c) hold. By
Lemma \ref{lemma-RS-implies-S1-S2}
we see that (S1) and (S2) hold. Hence by
Theorem \ref{theorem-miniversal-object-existence}
there exists a minimal versal formal object $\xi$ of $F$ such that
(\ref{equation-bijective-orbits}) holds. Say $\xi$ lies over $R$.
The map
$$
d\underline{\xi} : \text{Der}_\Lambda(R, k) \to T\mathcal{F}
$$
is bijective on $\text{Der}_\Lambda(k, k)$-orbits. Since the action
of $\text{Der}_\Lambda(k, k)$ on the left hand side is free by (c) and
Lemma \ref{lemma-action-linear}
we see that the map is bijective. Thus we see that $\underline{\xi}$
is an isomorphism by
Lemma \ref{lemma-minimal-smooth-morphism-functors}.
\end{proof}

\end{document}
